% !TeX root = ../papers/nc-p3.tex
\section{指数四における八次中心パラメータの存在と84次判定}
\label{sec:index-four-octic-normalization}
\begin{theorem}\label{thm:index-four-octic-normalization}
有限四辺の標準正則窓が中心上有限で頂点指数が4以下なら、経路長8の四つの斉次中心パラメータが存在する。
従って任意の成分でこの範囲は、八次の中心四元について中心剰余を29次、経路次元を84次、
Koszul複体を32次まで調べ、一般隅のPI次数が4以下であることを有限乗法表で確かめれば判定できる。
十五次元成分閉包の有限四辺では指数の追加検査は不要で、正則かつ一隅半素、正則かつ一隅整域、
正則かつ左右Noetherの三条件がいずれもこの84次判定に一致する。
\end{theorem}
\begin{proof}[四次元の単純における行列問題]
指数2以下はClifford、3は存在しない。指数4では定理\ref{thm:index-four-quartic-basepoints}により、
八次系に残りうる基点は全頂点四次元の単純 $S$ のみである。その矢空間を $H_i\subset\operatorname{Hom}(U_i,U_{i+1})$
と書く。各非零ベクトルで像の次元が2以上、各超平面の像の張る空間は全体である。双対でも同じである。
八次基点なら四次中心元の平方も消える。一般隅の被約traceはこの単純の行列traceに一致するので、
任意の始点で四段積空間 $L$ は $\operatorname{tr}L=0,\ \operatorname{tr}(LL)=0$ を満たすと仮定して矛盾を導く。
従って $L$ は十五次元の $\mathfrak{sl}_4$ の非退化trace形式で全等方であり $\dim L\le7$。
以下の行列論証はこの単純の存在と像の条件だけを用いる。

一辺でもどの二平面の像も三次元以上を張るなら、その辺を中央に置く三段積 $W$ は全非零ベクトル上で全射である。
最初の一段で二次元以上、次で三次元以上、最後で全体となる。続く辺を $H$ とすれば $HW=L$ もこの性質を持つ。
$p$ 次元から $q$ 次元への行列空間が全非零ベクトル上で全射なら、次元は $p+q-1$ 以上である。
双対零化空間が階数1のSegreと交わらないことによる。従って $\dim L=7$。
$H$ には非零な特異行列 $h$ がある（次元2以上の行列空間の行列式）。その階数を $r$ とすると、
$L$ の終域を $\operatorname{im} h$ で割った像も各ベクトル上全射なので次元は $7-r$ 以上、
従って $\operatorname{im} h$ に入る部分の次元は $r$ 以下である。だがそこに入る $hW$ は各非零ベクトル上で
$\operatorname{im}h$ に全射で次元 $3+r$ 以上である。矛盾する。
よって各 $H_i$ に、二平面 $P$ と二平面 $Q$ で $H_iP=Q$ となるものがある。
その部分と商の二行二列作用はともに各非零ベクトル上で全射である。商については $Q^\perp$ の非零共変元の
双対像が $P^\perp$ 全体であることを使う。従って両作用は行列空間として次元3または4である。
\end{proof}

\begin{proof}[階数二および階数一の矢の排除]
$h\in H_0$ が階数2であるとし、$P=\operatorname{im}h, K=\ker h$、$W=H_3H_2H_1$ と置く。
$hW$ を $P$ に制限した行列空間はtraceも二行列の積のtraceも零である。終域が $P$ に入るため
これは元の四段・八段のtraceと一致する。$\mathfrak{sl}_2$ の全等方空間なので次元は高々1で、
非零なら冪零行列の直線である。よってある $0\ne u\in P$ で $Wu\subset K$。
$H_1u=P_2$, $H_2 P_2=P_3$, $H_3 P_3=K$ は全て二平面である。途中で三次元に達すれば次段が全体となる。
双対でも同じことを逆向きに適用すると $H_1^*$ がある二平面を $P^\perp$ に送る。
従って $H_1(P)$ はその二平面の零化内に入り、実際に $H_1P=P_2$ である。

$H_1$ の $P\to P_2$ の作用には階数1の行列がある。全非零ベクトル上で全射な二行二列空間なので
次元3以上で、行列式二次曲面と交わるからである。これを与える $b$ に対し $bh=u_2\ell$ は階数1で
$u_2\in P_2$。八段traceの消滅から、続く六辺を $u_2$ に作用させた像は $\ker\ell$ に入る。
その最初の二段は $P_3,K$ 全体を与え、第三段 $H_0K$ は二平面でなければならない。
三次元以上ならその次に全体となるからである。双対版により $H_0^*(P^\perp)$ も二平面で、
従ってある二平面 $I$ に対し $H_0 I=P$ である。$I=K$ なら上の $P,P_2,P_3,K$ は巡回部分表現であり
単純性に反する。異なる二平面 $I,K$ は補空間になる。実際像が同じならその和に含まれる超平面の像が
全体にならず、像が異なれば共通する非零ベクトルの矢像の次元2以上に反する。同じ理由で両像も補空間である。
従って $H_0$ の矢は二つの二行二列作用の直和で書ける。

この直和において、両対角が\emph{ともに階数1}である矢 $h'$ が存在する。$\mathbb P(H_0)$ の次元が2なら
両対角への射は単射で、二つの射影二次曲線は交わる。次元3なら行列式二曲面の交叉は次元1以上で、
どちらかの対角が零になる高々二点を避けられる。$h'$ も階数2なので、先ほど示した必要条件から
$H_0\ker h'$ は二平面となるはずである。しかしその核は二つの元の平面内の直線の直和で、
各直線が対応する像平面全体を張るから $H_0\ker h'$ は四次元である。矛盾する。
以上はどの位置でも通る。また一般に、二つの異なる二平面で像がともに二平面なら同じ補空間の議論で
矢が直和になり、階数2が生じる。従って今は各辺のこうした入力二平面は一つだけである。

階数1の $h=u\ell\in H_0$ も存在しない。続く七段を $u$ に作用させた空間は $\ker\ell$ 内なので、
最初から第六段までの各像は二平面である。第一段と第五段の像は $H_2$ が二平面へ送るので、
その唯一性から同じである。その間の四辺が二平面の巡回部分表現を与え、再び矛盾する。
\end{proof}

\begin{proof}[唯一の二平面と特異矢の像]
一辺 $H$ をその二平面の部分と商について上三角に書き、二つの対角矢を $A,C$、右上を $D$ とする。
非零特異矢は今や全て階数3である。$\dim H=4$ なら $A,C$ はどちらも $H\to M_2$ の同型である。
どちらかの核の非零矢は全階数が2以下になってしまう。二つの行列式は滑らかな二次曲面を定める。
異なればその完全交叉 $\Gamma$ を、同じならその曲面上の一般の型 $(2,2)$ の曲線 $\Gamma$ を取る。
$\Gamma$ 上で $A,C:2\mathcal O(-1)\to2\mathcal O$ は階数一定1である。
$D$ の誘導する直線束の射 $\ker C\to\operatorname{coker} A$ が点で零になることは全矢の階数が2になることと同値である。
しかしこの射がどこでも可逆であることはできない。両余核はそれぞれの滑らかな二次曲面の一方の罫線束なので
$\Gamma$ で次数2である。$\mathcal O_\Gamma(1)$ の次数4と核・余核の行列式の式から $\deg\ker C=2-8=-6$。
次数が異なる。この計算は二曲面の交叉が非被約でも直線束の次数で成立する。

残る $\dim H=3$ では両対角の行列式が $\mathbb P(H)$ の滑らかな二次曲線を定める。
同じ曲線なら前の議論をその二次曲線上で繰り返すと余核次数1、核次数 $1-4=-3$ となって矛盾する。
従って相異なる。全ての特異矢の像超平面の交わりは零である。実際 $C$ の行列式曲線で $A$ が可逆な部分では
像超平面を対角の出力二平面 $Q$ で割った直線が動き、交わりは $Q$ 内に入る。
逆に $A$ の行列式曲線で $C$ が可逆な部分では像超平面と $Q$ の交わりが $\operatorname{im} A$ という動く直線である。
どちらの曲線でも階数1行列の像直線への射は同型であり、有限個の交点を外しても一定にはならない。

各辺の唯一の二平面の出力が次辺の入力に常に一致すれば、これらが部分表現となる。従って例えば
$H_1$ の出力二平面と $H_2$ の入力二平面 $P_2$ が異なるように巡回してよい。
$W=H_3H_2H_1$ が $v\in U_1$ 上で全射でなければ、$H_1v$ は二平面で $H_2$ はそれを二平面に送るので
$H_1v=P_2$ となる。そのような $v$ は線型空間 $\{v:H_1v\subset P_2\}$ に属し、その次元は高々1である。
二次元あれば $H_1$ の唯一の出力が $P_2$ となり、三次元あれば超平面の全射性に反する。
この例外直線を含まない像超平面 $J$ を持つ特異矢 $h\in H_0$ が前段により存在する。
$W$ は $J$ の全非零ベクトル上で全射なので、$(hW)|_J$ は三次元間で各非零ベクトル上全射、従ってその次元は5以上である。
一方そのtraceおよび二要素の積のtraceは元の四段・八段で計算され零である。
八次元の $\mathfrak{sl}_3$ の全等方空間の次元は4以下であり、最後の矛盾を得る。

これで八次中心系の基点はなく、三次元射影中心をその四一般元で切って有限剰余を得る。
有限検査は定理\ref{thm:successively-central-cutoff}に $d_i=8$ を入れたもの。
逆に検査を通った場合の中心有限性と素性は従前の中心切断の証明による。指定閉包は常に十六行十六列行列の
恒等式を満たすので、正則整域・半素・Noetherの既出の同値と今回の必要性で最後の判定が従う。
\end{proof}
